Below is a LaTeX-ready writeup of \((1)\) and \((2)\).  
For \((3)\), the excerpt you gave does **not** define the measures \(\mu^\pm\), so a completely rigorous proof of the statement about
\[
w=\frac{d(\mu^++\mu^-)}{dA_{\mathrm{sph}}}
\]
is impossible without guessing what \(\mu^\pm\) are. In keeping with your instruction not to guess, I will prove \((1)\) and \((2)\) completely, and then state the exact positivity/symmetry facts that follow from them and are typically what one uses to finish \((3)\) once \(\mu^\pm\) are specified.

---

Let
\[
a=\sqrt{\frac23},\qquad b=\frac1{\sqrt3},
\]
and let
\[
q=(x,y,z)\in \mathbb S^2\setminus\{p_1,\dots,p_8\}.
\]
Write the Green function in the form
\[
G_p(q)=-\frac1{4\pi}\log(1-p\cdot q)+c_0,
\]
where \(c_0\) is the fixed additive constant coming from the chosen normalization.

Since both \(p\) and \(q\) are unit vectors and \(q\neq p\), one has \(p\cdot q<1\), hence
\[
1-p\cdot q>0.
\]
Therefore all logarithms below are well-defined.

## Proof of \((1)\)

For the \(A\)-vertices we have
\[
p_1=(b,0,a),\quad p_2=(-b,0,a),\quad p_3=(b,0,-a),\quad p_4=(-b,0,-a).
\]
Hence
\[
p_1\cdot q=bx+az,\quad
p_2\cdot q=-bx+az,\quad
p_3\cdot q=bx-az,\quad
p_4\cdot q=-bx-az.
\]
Therefore
\[
\begin{aligned}
G_A(q)
&=\sum_{i=1}^4 G_{p_i}(q) \\
&=-\frac1{4\pi}\sum_{i=1}^4 \log(1-p_i\cdot q)+4c_0 \\
&=-\frac1{4\pi}\log\!\Bigl(\prod_{i=1}^4 (1-p_i\cdot q)\Bigr)+C_A,
\end{aligned}
\]
where \(C_A=4c_0\).

Now
\[
\begin{aligned}
\prod_{i=1}^4 (1-p_i\cdot q)
&=(1-bx-az)(1+bx-az)(1-bx+az)(1+bx+az) \\
&=\bigl((1-az)^2-b^2x^2\bigr)\bigl((1+az)^2-b^2x^2\bigr) \\
&=P_A(x,z).
\end{aligned}
\]
Thus
\[
G_A(q)=-\frac1{4\pi}\log P_A(x,z)+C_A.
\]

The \(B\)-vertices are
\[
p_5=(b,a,0),\quad p_6=(-b,a,0),\quad p_7=(-b,-a,0),\quad p_8=(b,-a,0),
\]
so
\[
p_5\cdot q=bx+ay,\quad
p_6\cdot q=-bx+ay,\quad
p_7\cdot q=-bx-ay,\quad
p_8\cdot q=bx-ay.
\]
Exactly as above,
\[
\begin{aligned}
G_B(q)
&=\sum_{j=5}^8 G_{p_j}(q) \\
&=-\frac1{4\pi}\log\!\Bigl(\prod_{j=5}^8 (1-p_j\cdot q)\Bigr)+C_B,
\end{aligned}
\]
with \(C_B=4c_0\), and
\[
\begin{aligned}
\prod_{j=5}^8 (1-p_j\cdot q)
&=(1-bx-ay)(1+bx-ay)(1+bx+ay)(1-bx+ay) \\
&=\bigl((1-ay)^2-b^2x^2\bigr)\bigl((1+ay)^2-b^2x^2\bigr) \\
&=P_B(x,y).
\end{aligned}
\]
Hence
\[
G_B(q)=-\frac1{4\pi}\log P_B(x,y)+C_B.
\]

This proves \((1)\).

---

## Proof of \((2)\)

### The derivative of \(G_A\)

Set
\[
U_A=(1-az)^2-b^2x^2,\qquad V_A=(1+az)^2-b^2x^2,
\]
so that
\[
P_A=U_A V_A.
\]
Since \(G_A\) depends only on \(x\) and \(z\), we have
\[
\partial_y G_A=0.
\]
Also
\[
\partial_x U_A=-2b^2x,\qquad \partial_x V_A=-2b^2x,
\]
hence
\[
\begin{aligned}
\partial_x P_A
&=(\partial_x U_A)V_A+U_A(\partial_x V_A) \\
&=-2b^2x(U_A+V_A).
\end{aligned}
\]
Now
\[
U_A+V_A
=(1-az)^2+(1+az)^2-2b^2x^2
=2\bigl(1+a^2z^2-b^2x^2\bigr),
\]
so
\[
\partial_x P_A=-4b^2x\bigl(1+a^2z^2-b^2x^2\bigr).
\]
Therefore
\[
\begin{aligned}
\partial_x G_A
&=-\frac1{4\pi}\frac{\partial_x P_A}{P_A} \\
&=\frac{b^2x\bigl(1+a^2z^2-b^2x^2\bigr)}{\pi P_A}.
\end{aligned}
\]
Using
\[
a^2=\frac23,\qquad b^2=\frac13,
\]
we obtain
\[
b^2\bigl(1+a^2z^2-b^2x^2\bigr)
=\frac13\left(1+\frac23 z^2-\frac13 x^2\right)
=\frac{3+2z^2-x^2}{9}.
\]
Thus
\[
\partial_x G_A(q)=\frac{x(3+2z^2-x^2)}{9\pi\,P_A(x,z)}.
\]
Since
\[
\partial_\phi=-y\,\partial_x+x\,\partial_y,
\]
and \(\partial_y G_A=0\), we get
\[
\partial_\phi G_A(q)
=-y\,\partial_x G_A(q)
=-\frac{xy(3+2z^2-x^2)}{9\pi\,P_A(x,z)}.
\]

### The derivative of \(G_B\)

Write
\[
U_B=(1-ay)^2-b^2x^2,\qquad V_B=(1+ay)^2-b^2x^2,
\]
so that
\[
P_B=U_BV_B.
\]
It is convenient to note that
\[
P_B=(1+a^2y^2-b^2x^2)^2-4a^2y^2.
\]

First,
\[
\partial_x P_B=-4b^2x\bigl(1+a^2y^2-b^2x^2\bigr),
\]
so
\[
\begin{aligned}
\partial_x G_B
&=-\frac1{4\pi}\frac{\partial_x P_B}{P_B} \\
&=\frac{b^2x\bigl(1+a^2y^2-b^2x^2\bigr)}{\pi P_B}.
\end{aligned}
\]
Again using \(a^2=\frac23\) and \(b^2=\frac13\),
\[
b^2\bigl(1+a^2y^2-b^2x^2\bigr)
=\frac13\left(1+\frac23 y^2-\frac13 x^2\right)
=\frac{3+2y^2-x^2}{9},
\]
hence
\[
\partial_x G_B(q)=\frac{x(3+2y^2-x^2)}{9\pi\,P_B(x,y)}.
\]

Next,
\[
\begin{aligned}
\partial_y P_B
&=2(1+a^2y^2-b^2x^2)(2a^2y)-8a^2y \\
&=4a^2y\bigl(a^2y^2-b^2x^2-1\bigr) \\
&=-4a^2y\bigl(1-a^2y^2+b^2x^2\bigr).
\end{aligned}
\]
Therefore
\[
\begin{aligned}
\partial_y G_B
&=-\frac1{4\pi}\frac{\partial_y P_B}{P_B} \\
&=\frac{a^2y\bigl(1-a^2y^2+b^2x^2\bigr)}{\pi P_B}.
\end{aligned}
\]
Substituting \(a^2=\frac23\) and \(b^2=\frac13\),
\[
a^2\bigl(1-a^2y^2+b^2x^2\bigr)
=\frac23\left(1-\frac23 y^2+\frac13 x^2\right)
=\frac{2(3+x^2-2y^2)}{9},
\]
so
\[
\partial_y G_B(q)=\frac{2y(3+x^2-2y^2)}{9\pi\,P_B(x,y)}.
\]

Now apply \(\partial_\phi=-y\,\partial_x+x\,\partial_y\):
\[
\begin{aligned}
\partial_\phi G_B(q)
&=-y\,\partial_x G_B(q)+x\,\partial_y G_B(q) \\
&=\frac{xy}{9\pi\,P_B(x,y)}
\Bigl(-(3+2y^2-x^2)+2(3+x^2-2y^2)\Bigr) \\
&=\frac{xy}{9\pi\,P_B(x,y)}\bigl(3+3x^2-6y^2\bigr) \\
&=\frac{xy(1+x^2-2y^2)}{3\pi\,P_B(x,y)}.
\end{aligned}
\]

This proves \((2)\).

---

## What can be proved rigorously toward \((3)\) from the supplied data

As written, \((3)\) is incomplete because \(\mu^\pm\) are not defined. Still, the key facts needed for that step are immediate.

Define
\[
R_x(x,y,z)=(x,z,-y).
\]
Then
\[
\begin{aligned}
P_A(R_x q)
&=((1-a(-y))^2-b^2x^2)\,((1+a(-y))^2-b^2x^2) \\
&=((1+ay)^2-b^2x^2)\,((1-ay)^2-b^2x^2) \\
&=P_B(q),
\end{aligned}
\]
and similarly
\[
P_B(R_x q)=P_A(q).
\]
Hence \(R_x\) swaps the \(A\)- and \(B\)-polynomials.

Also,
\[
P_A(q)=\prod_{i=1}^4 (1-p_i\cdot q),\qquad
P_B(q)=\prod_{j=5}^8 (1-p_j\cdot q),
\]
and each factor is strictly positive on
\[
\mathbb S^2\setminus\{p_1,\dots,p_8\}.
\]
Therefore
\[
P_A(q)>0,\qquad P_B(q)>0
\quad\text{for all }q\in \mathbb S^2\setminus\{p_1,\dots,p_8\}.
\]

Moreover, with the same Green-function normalization for every \(G_p\), one has \(C_A=C_B\), and so
\[
G_A(R_x q)=G_B(q),\qquad G_B(R_x q)=G_A(q).
\]

These are exactly the symmetry/positivity inputs one normally uses to prove that a density \(w\) built symmetrically from the \(A\)- and \(B\)-contributions is strictly positive and \(R_x\)-invariant. But to finish the literal statement of \((3)\), I would need the explicit definition of \(\mu^\pm\).

If you send that definition, I can complete \((3)\) in the same style.